% Chapter 3 -- Methodology: problem formulation, the placement taxonomy, the four methods,
% the offline re-scoring procedures, the post-tuning experiment, and the threat model.
% Status: basic draft; formulas follow the module code. Pending: attack protocol details.
\chapter{Methodology}
\label{ch:method}

\section{Problem Formulation}
\label{sec:method-problem}

Let $\mathcal{U}$ be the set of users and $\mathcal{I}$ the set of items. Each user $u$ has an
ordered interaction history; the latest interaction $t_u$ is held out as the test item and the
remaining interactions form the local training set $\mathcal{D}_u$, all treated as implicit
positives. In the cross-device protocol each user is one client, so the client population is
$\mathcal{U}$ with $N=|\mathcal{U}|=\NClients$, and the federated algorithm runs for
$R=\NRounds$ rounds, sampling $\lceil C N\rceil = \ClientsPerRound$ clients per round with
$C=\ParticipationC$. A client therefore participates about $CR = \ExpectedParticipations$ times
in expectation and keeps its private parameters on the device between participations.

Every method scores a user--item pair with a function $\hat{s}_{ui} = f(\theta_u, \vect{w};
u, i)$ of private parameters $\theta_u$ and global parameters $\vect{w}$. The \emph{placement}
of personalisation is the choice of what goes into $\theta_u$. At evaluation the candidate set
$\mathcal{C}_u$ of user $u$ is $t_u$ plus 99 items sampled from the items $u$ never interacted
with, identical across methods (Section~\ref{sec:setup-protocol}); the metrics are \hr{},
\ndcg{} and \mrr{} of $t_u$ within $\mathcal{C}_u$, averaged over all $N$ users and within the
sparse, medium and dense groups.

\section{A Taxonomy of Placements}
\label{sec:method-taxonomy}

Figure~\ref{fig:placement} shows the four placements compared in this thesis.

\begin{figure}[htbp]
  \centering
  % TikZ: the four placements of the personal part. Input from the methodology chapter.
\begin{tikzpicture}[
  font=\small,
  box/.style={draw, rounded corners=2pt, minimum width=2.75cm, minimum height=0.6cm, align=center},
  glob/.style={box, fill=blue!8},
  priv/.style={box, fill=red!8, densely dashed},
  srv/.style={draw, fill=gray!12, minimum width=2.75cm, minimum height=0.55cm, align=center},
  up/.style={-{Stealth[length=2mm]}, thick},
  lab/.style={font=\footnotesize\itshape, align=center}
]
  % column x positions
  \def\dx{3.3}
  \foreach \k/\name in {0/Baseline, 1/Personalized, 2/PFedRec, 3/Adaptive} {
    \node[srv] (srv\k) at (\k*\dx, 3.2) {server\\aggregates};
    \node[lab] at (\k*\dx, 4.0) {\textbf{\name}};
  }
  % Baseline: everything global
  \node[glob] (b1) at (0, 1.8) {$\pu$, $\qi$, biases};
  \node[lab]  at (0, 0.9) {no private part};
  \draw[up] (b1) -- (srv0) node[midway, right, font=\scriptsize] {all};
  % Personalized: user vector private
  \node[glob] (p1) at (\dx, 1.8) {$\qi$, biases};
  \node[priv] (p2) at (\dx, 0.9) {$\pu$ (on device)};
  \draw[up] (p1) -- (srv1) node[midway, right, font=\scriptsize] {globals};
  % PFedRec: private head
  \node[glob] (f1) at (2*\dx, 1.8) {$\Qtab$, $b$};
  \node[priv] (f2) at (2*\dx, 0.9) {$\wu$ (head)};
  \draw[up] (f1) -- (srv2) node[midway, right, font=\scriptsize] {item table};
  % Adaptive: blend + MLP, uploads p_u as prototype contribution
  \node[glob] (a1) at (3*\dx, 1.8) {$\qi$, biases, $\bar{\vect{p}}$};
  \node[priv] (a2) at (3*\dx, 0.9) {$\pu$, PersonalMLP, $\alpha_u$};
  \draw[up] (a1) -- (srv3) node[midway, right, font=\scriptsize] {globals};
  \draw[up, red!70!black] (a2.east) .. controls (3*\dx+2.1, 1.5) and (3*\dx+2.1, 2.8) .. (srv3.east)
    node[pos=0.5, right, font=\scriptsize, text=red!70!black] {$\pu$};
  % legend
  \node[glob, minimum width=1.6cm, minimum height=0.4cm] (lg) at (0.6, -0.3) {};
  \node[right=2pt of lg, font=\scriptsize] {global, aggregated};
  \node[priv, minimum width=1.6cm, minimum height=0.4cm] (lp) at (5.2, -0.3) {};
  \node[right=2pt of lp, font=\scriptsize] {private, stays on the device};
\end{tikzpicture}

  \caption{The four placements of the personal part. Blue boxes are aggregated by the server;
  dashed red boxes stay on the device. The adaptive method also uploads the user vector
  $\pu$ as its contribution to the global prototype (red arrow), which Section~\ref{sec:method-threat}
  treats as a leakage channel.}
  \label{fig:placement}
\end{figure}

\begin{enumerate}[leftmargin=2em]
  \item \textbf{None (all-global).} Every parameter, including the user vectors, is aggregated.
        A user vector is updated only when its owner participates, and the aggregation mixes it
        with the stale copies held by the other sampled clients.
  \item \textbf{Private user vector (split learning).} The user vector and bias never leave the
        device; the item table and biases are aggregated.
  \item \textbf{Private scoring head.} There is no user vector. Each client owns a linear head
        over the shared item table; the item table is aggregated.
  \item \textbf{Blend with population knowledge.} The private user vector is interpolated with a
        global prototype by a per-user weight $\alpha_u$, and a private neural scorer is added.
\end{enumerate}

\section{The Four Methods}
\label{sec:method-methods}

All four methods use the same data split, the same exclusion table for negative sampling and the
same scoring routine (Chapter~\ref{ch:setup}); they differ in the model, the loss and the
placement. Table~\ref{tab:hyperparams} in Appendix~\ref{app:config} lists the configurations of
the cited runs.

\subsection{Baseline: all-global BPR-MF}
\label{sec:method-baseline}

The model is matrix factorisation with biases,
\begin{equation}
  \hat{s}_{ui} = b + b_u + b_i + \pu^\top \qi, \qquad \pu, \qi \in \mathbb{R}^{d},
  \label{eq:mf}
\end{equation}
trained with the BPR loss \cite{rendle2009bpr}
\begin{equation}
  \mathcal{L}_u = -\sum_{i \in \mathcal{D}_u}\sum_{j \in \mathcal{N}_u(i)}
    \log \sigma\big(\hat{s}_{ui} - \hat{s}_{uj}\big) + \lambda \|\theta\|^2,
  \label{eq:bpr}
\end{equation}
where $\mathcal{N}_u(i)$ are negatives sampled from the items $u$ has not interacted with.
Every parameter is global: the server runs FedAvg \cite{mcmahan2017fedavg}, weighting the
sampled clients by their number of positives. The client uploads \ParamsBaselineGlobal{}
parameters per participation. This placement is the lower bound: its user vectors are averaged
with copies that other clients never trained.

\subsection{Personalized: split learning with a private user vector}
\label{sec:method-personalized}

The same model and loss as the baseline, with the parameter set split into
$\theta_u = \{\pu, b_u\}$ (private, cached on the device) and $\vect{w} = \{\Qtab, \vect{b}_I,
b\}$ (global). In each participation the client loads $\theta_u$ from its cache, trains all
parameters on $\mathcal{D}_u$, saves $\theta_u$ and returns only $\vect{w}$. The FedProx variant
regularises the global part only, $\frac{\mu}{2}\|\vect{w} - \vect{w}_{\text{server}}\|^2$, so
that the private part is free to diverge. The upload is \ParamsSplitGlobal{} parameters. This
is the placement of FCF and FedMF \cite{ammaduddin2019fcf,chai2021fedmf}, and the
``knob-for-knob'' partner of the baseline: the two cited runs share the loss, the dimension
$d=\DimMf$, the optimiser and the learning rate, and differ only in whether $\pu$ is averaged.

\subsection{PFedRec: private scoring head on a shared item table}
\label{sec:method-pfedrec}

PFedRec \cite{zhang2023pfedrec} has no user vector. Each client owns a linear head
$\wu \in \mathbb{R}^{d}$ and scores
\begin{equation}
  \hat{s}_{ui} = \sigma\big(\wu^\top \qi + b\big), \qquad d=\DimPfedrec,
  \label{eq:pfedrec}
\end{equation}
with a BCE loss on the observed items (positives) and four sampled negatives per positive. One
local epoch alternates, per mini-batch, a step on the head with learning rate $\eta_w$ and a
step on the item rows with learning rate $\eta_w \cdot |\mathcal{I}| \cdot \eta$; the large item
learning rate compensates for the sparsity of the item gradients. The server averages the item
tables uniformly over the participating clients, as in the reference implementation. The upload
is \ParamsPfedrecGlobal{} parameters, the smallest of the four.

Two evaluation details matter under $C<1$. First, a client that has not participated for many
rounds holds a head that was trained against an older item table. The cited run therefore ends
with a \emph{final calibration}: after the best-round item table is restored, every client runs
one local epoch on the head only, with the item table frozen, so that each head is aligned to
the table that is evaluated. Second, the reference implementation scores each user with the
client's own post-tuned item table $\Qu$ from its last local epoch; the thesis number scores
with the shared table $\Qtab$, and Section~\ref{sec:method-posttune} measures the difference.

\subsection{Adaptive: dual-level personalisation with a hierarchical-conditional $\alpha$}
\label{sec:method-adaptive}

The adaptive method is the thesis's own design on top of the split-learning placement. It has
two levels.

\paragraph{Level 1: blended user vector.} The server keeps a global prototype
$\bar{\vect{p}}$, an exponential moving average with momentum \PrototypeMomentum{} of the
user vectors uploaded by the participating clients, weighted by their number of positives. The
client forms
\begin{equation}
  \tilde{\vect{p}}_u = \alpha_u \pu + (1-\alpha_u)\, \bar{\vect{p}},
  \label{eq:blend}
\end{equation}
where $\alpha_u \in [\AlphaMin, \AlphaMax]$ is computed from four per-user statistics (interaction count,
genre entropy, number of unique items, rating standard deviation) by a two-stage rule: a
hierarchical aggregation of a data-volume score and a preference-quality score, followed by
conditional adjustments for sparse users (down to a 50\% penalty), niche specialists ($+0.15$),
inconsistent raters (30\% penalty) and completionists ($+0.10$). The rule is deterministic and
needs no gradient; its ablation partners are a naive multi-factor weighted sum and a single
data-quantity factor.

\paragraph{Level 2: private neural scorer.} A per-client multi-layer perceptron scores the
element-wise product of the blended user vector and the item vector,
\begin{equation}
  \hat{s}^{\text{MLP}}_{ui} = \text{PersonalMLP}_u\big(\tilde{\vect{p}}_u \odot \qi\big),
  \qquad
  \hat{s}_{ui} = \vect{v}_u^\top \big[\hat{s}^{\text{CF}}_{ui},\ \hat{s}^{\text{MLP}}_{ui}\big] + c_u,
  \label{eq:dual}
\end{equation}
where $\hat{s}^{\text{CF}}_{ui} = b + b_u + b_i + \tilde{\vect{p}}_u^\top\qi$ and the fusion
$(\vect{v}_u, c_u)$ is a private linear layer. With hidden sizes \MlpHidden{} the PersonalMLP
has \ParamsPersonalMlp{} private parameters, against a median of \DsMedianRatings{}
interactions per user. Training uses the BPR loss, FedProx on the global part, and a contrastive
term between $\pu$ and $\tilde{\vect{p}}_u$ with weight \ContrastiveLambda; the optional
per-user learnable $\alpha_u$ and item perturbation were disabled in the cited run.

\paragraph{What is transmitted.} The global part (\ParamsSplitGlobal{} parameters) and, in
addition, the user vector $\pu$ (\ParamsAdaptiveUploadExtra{} floats) as the client's
contribution to the prototype. The second upload is what distinguishes this placement from
split learning in the threat model.

\section{Offline Re-scoring and Counterfactual Heads}
\label{sec:method-reeval}

To compare scoring conditions on the same trained model, an offline re-scoring tool reproduces
the final federated evaluation: it loads the best-round global parameters of a run, the cached
private parameters of every client, and the same candidate list per user, and recomputes the
metrics. A \emph{replay trust gate} requires that this reproduction match the run's recorded
per-user and aggregate metrics exactly before any variant is scored; on the cited PFedRec run the
match is bit-exact (Section~\ref{sec:setup-validity}).

On the PFedRec model the tool then answers the placement question with counterfactual heads,
holding $\Qtab$, $b$ and the candidate lists fixed and replacing $\wu$ at scoring time by:
\begin{itemize}[leftmargin=2em]
  \item the population-mean head $\bar{\vect{w}}$ (uniform mean of all $N$ calibrated heads),
        and a variant weighted by interaction count;
  \item the blend $\alpha_u \wu + (1-\alpha_u)\bar{\vect{w}}$ for every $\alpha$ on a grid of
        \HeadAlphaGridSize{} values with step \HeadAlphaGridStep, and for the
        hierarchical-conditional rule of Section~\ref{sec:method-adaptive} applied to the
        PFedRec head;
  \item another user's head (a random permutation of the heads);
  \item a freshly initialised head per user, and one shared freshly initialised head, which
        reproduce the cold-evaluation failure deliberately.
\end{itemize}
Each variant yields a per-user \ndcg{}, so the comparison with the own-head reference is paired.

\section{Post-tuning Re-score of PFedRec}
\label{sec:method-posttune}

The second re-scoring experiment measures how much of the gap between the reproduced PFedRec
number and the published one is due to the reference implementation's inference protocol. It
was pre-registered (Appendix~\ref{app:posttune}) and then run once. For every user $u$:
\begin{enumerate}[leftmargin=2em]
  \item initialise $\Qu \leftarrow \Qtab$ (the restored best-round table) and freeze the
        calibrated head $(\wu, b)$;
  \item fine-tune only the rows of $\Qu$ on $\mathcal{D}_u$ for $E \in \{1, 3\}$ epochs with the
        run's own item learning rate (\PtLrItem), BCE, four negatives per positive, batch 256,
        plain SGD;
  \item score $u$ with $\Qu$ on the same 100 candidates.
\end{enumerate}
Step 2 is run with two negative-sampling pools. The \emph{clean} pool draws negatives from all
items $u$ has not trained on, so the held-out test item can be drawn, as a client that does not
know the future would do. The \emph{official-code} pool additionally excludes the test item, as
the released PFedRec code does. The two pools are the only difference between the two
trajectories, so the contrast official $-$ clean isolates the test-item leak, and the contrast
clean $-$ reference isolates the effect of post-tuning itself. Table~\ref{tab:posttune-deviations}
lists the known deviations from the reference implementation.

\begin{table}[htbp]
  \caption{Known deviations of the post-tuning re-score from the reference implementation and
  their expected effect (registered before the run).}
  \label{tab:posttune-deviations}
  \centering\small
  \begin{tabular}{p{0.05\linewidth}p{0.3\linewidth}p{0.28\linewidth}p{0.3\linewidth}}
    \toprule
    \# & Reference code & This experiment & Consequence \\
    \midrule
    K1 & Alternates head and item updates per batch & Head frozen, item rows only & Post-tuning is weaker here; the head was already calibrated \\
    K2 & Private table is the one after the user's last participation ($C=1$: always the last round) & Every user starts from the restored best-round table & Removes staleness from the measurement; isolates post-tuning \\
    K3 & Validation item selects the round & None; best round selected on the test item & Inflates our numbers (declared limitation) \\
    K4 & Test item excluded from the negatives & Both pools run & Separates the leak from post-tuning \\
    K5 & $C=1$ & $C=\ParticipationC$, unchanged & The residual gap is left as ``not separated'' \\
    \bottomrule
  \end{tabular}
\end{table}

\section{Threat Model and Privacy Measurements}
\label{sec:method-threat}

\paragraph{Adversary.} The server is \emph{honest-but-curious}: it follows the protocol, stores
every message it receives, and tries to infer the users' interactions and preferences. Clients
are honest. No secure aggregation is deployed in the cited runs, so the server sees every
client's upload individually.

\paragraph{Leakage channels per placement.}
\begin{itemize}[leftmargin=2em]
  \item \emph{All-global:} the full user vector and bias of the participating user, and the
        update pattern of the item table (non-zero rows = interacted items plus sampled
        negatives).
  \item \emph{Private user vector:} the item-table update pattern only.
  \item \emph{Private head:} the item-table update pattern only; the head never leaves the device.
  \item \emph{Blend with population knowledge:} the item-table update pattern \emph{and} the
        user vector $\pu$, uploaded as the prototype contribution.
\end{itemize}
The update pattern is common to all placements and is mitigated by secure aggregation
\cite{bonawitz2017secagg}, which reveals only the sum over the sampled clients. The upload of
$\pu$ is specific to the adaptive placement and is not mitigated by secure aggregation of the
item table.

\paragraph{Measurements.} Two quantities are reported in Section~\ref{sec:res-privacy}.
(i) The communication volume per participating client per round, counted in parameters from
the architectures (Table~\ref{tab:params-comm}) and \pending{measured in bytes from the actual
Flower messages}. (ii) A \emph{preference-reconstruction attack}: for every user of the adaptive
run, the server ranks all $|\mathcal{I}|$ items by $\pu^\top \qi$ with the received $\pu$ and
its own item table, takes the top 10, and counts how many of them the user actually interacted
with; the score is compared with a random ranking and with a popularity ranking.
\pending{attack protocol to be pre-registered and run; report hit counts and the two baselines
overall and by group, and the effect of secure aggregation as the mitigation}.
